\documentclass[12pt,a4paper]{article}
\usepackage{ugmath}
\usepackage{float}
\usepackage{placeins}
\usepackage[labelfont=bf,labelsep=space]{caption}
\newcommand{\paren}[1]{\left( #1 \right)}
\newcommand{\alumno}{Ricardo León Martínez}
\newcommand{\materia}{Fisica I}
\newcommand{\profesor}{Lucero Uscanga Aguilera}
\newcommand{\tarea}{Tarea 2}
\newcommand{\fecha}{24/2/2026}

\begin{document}
    \begin{center}
    {\large\textbf{UNIVERSIDAD DE GUANAJUATO}}\\[0.3cm]
    {\normalsize\textbf{DIVISIÓN DE CIENCIAS NATURALES Y EXACTAS}}\\
    {\normalsize\textbf{CAMPUS GUANAJUATO}}\\[1cm]

    {\Large\textbf{\tarea\ (\materia)}}\\[1cm]
    \end{center}

    \textbf{Nombre:} \alumno \hfill 
    \textbf{Fecha:} \fecha \hfill 
    \textbf{Calificación:} \rule{3cm}{0.4pt} \\[0.3cm]

    \begin{mdframed}[style=mdbluebox,frametitle={Ejercicio 1}]
        La viga $AB$ es uniforme y tiene una masa de $100kg$. Descansa en sus extremos $A$ y $B$
        y soporta las masas como se indica en la figura. Calcular la reacción en los soportes.
    \end{mdframed}

    \begin{solucion}
        Denotemos por $R_A$ y $R_B$ las reacciones en los apoyos $A$ y $B$.
        Como el sistema se encuentra en equilibrio estático, se satisfacen las
        ecuaciones
        \begin{equation*}
            \sum F_y = 0
            \qquad
            \sum \tau_A = 0
        \end{equation*}
        La viga es uniforme, por tanto su peso actúa en su centro geométrico,
        esto es, a una distancia de $2\,\text{m}$ del punto $A$. Sea $g=9.8\,\text{m/s}^2$. Entonces los pesos correspondientes son
        \begin{align*}
            P_1 &= 50g, \\
            P_2 &= 150g, \\
            P_v &= 100g.
        \end{align*}
        La ecuación de equilibrio de fuerzas verticales en este caso es
        \begin{equation*}
            R_A + R_B - P_1 - P_2 - P_v = 0
        \end{equation*}
        Sustituyendo,
        \begin{equation*}
            R_A + R_B - 50g - 150g - 100g = 0
        \end{equation*}
        \begin{equation*}
            R_A + R_B = 300g
        \end{equation*}
        Tomando momentos positivos en sentido antihorario,

        \begin{equation*}
            R_B(4) - P_1(1) - P_v(2) - P_2(3) = 0
        \end{equation*}

        Sustituyendo los valores,

        \begin{equation*}
            4R_B - 50g(1) - 100g(2) - 150g(3) = 0
        \end{equation*}

        \begin{equation*}
            4R_B - 50g - 200g - 450g = 0
        \end{equation*}

        \begin{equation*}
            4R_B - 700g = 0
        \end{equation*}

        \begin{equation*}
            R_B = 175g
        \end{equation*}

        Sustituyendo en la ecuación de fuerzas,

        \begin{equation*}
            R_A + 175g = 300g
        \end{equation*}

        \begin{equation*}
            R_A = 125g
        \end{equation*}

        \textbf{Conclusión}

        Las reacciones en los apoyos son

        \begin{align*}
            R_A &= 125g = 1225\,\text{N}, \\
            R_B &= 175g = 1715\,\text{N}.
        \end{align*}
    \end{solucion}

    \begin{mdframed}[style=mdbluebox,frametitle={Ejercicio 2}]
        La viga uniforme $AB$ de la figura tiene $4cm$ de largo y pesa $100kgf$. La viga puede
        rotar alrededor del punto fijo $C$. La viga reposa en el punto $A$. Un hombre que pesa
        $75kgf$ camina a lo largo de la viga, partiendo de $A$. Calcular la máxima distancia que
        el hombre puede caminar a partir de $A$ manteniendo el equilibrio. Representar
        la fuerza de reaccion en $A$ como una función de la distancia $x$.
    \end{mdframed}

    \begin{solucion}
        Sea $AB$ una viga uniforme de longitud $4\,\text{m}$ y peso 
        $P_v=100\,\text{kgf}$.  

        Los puntos están dispuestos en el orden $A-C-B$ con

        \begin{equation*}
            AC = 2.5\,\text{m}.
        \end{equation*}

        El punto $C$ es el pivote fijo.
        En $A$ existe una reacción vertical $R_A$.
        Un hombre de peso $P_h=75\,\text{kgf}$ se encuentra a una distancia $x$
        medida desde $A$. El centro de masa de la viga está a $2\text{ m}$
        desde $A$. Tomamos momentos positivos en sentido antihorario.
        La distancia de $A$ a $C$ es $2.5\,\text{m}$. La distancia del centro de masa al punto $C$ es
        \begin{equation*}
            2 - 2.5 = -0.5\,\text{m}.
        \end{equation*}
        La distancia del hombre al punto $C$ es
        \begin{equation*}
            x - 2.5.
        \end{equation*}
        La ecuación de momentos queda
        \begin{equation*}
            - R_A(2.5) + 100(-0.5) + 75(x-2.5) = 0
        \end{equation*}
        \begin{equation*}
            -2.5R_A -50 + 75x -187.5 = 0
        \end{equation*}
        \begin{equation*}
            -2.5R_A + 75x -237.5 = 0
        \end{equation*}
        \begin{equation*}
            2.5R_A = 75x -237.5
        \end{equation*}
        \begin{equation*}
            R_A(x) = \frac{75x -237.5}{2.5}
        \end{equation*}
        El equilibrio requiere
        \begin{equation*}
            R_A(x) \ge 0.
        \end{equation*}
        El valor crítico ocurre cuando
        \begin{equation*}
            R_A(x)=0.
        \end{equation*}
        Entonces
        \begin{equation*}
            75x -237.5 = 0
        \end{equation*}
        \begin{equation*}
            x = \frac{237.5}{75}
        \end{equation*}
        \begin{equation*}
            x = \frac{19}{6}
        \end{equation*}
        \begin{equation*}
            x \approx 3.17\,\text{m}.
        \end{equation*}
        La reacción en $A$ es
        \begin{equation*}
            R_A(x)=\frac{75x-237.5}{2.5}
        \end{equation*}
        y la distancia máxima que el hombre puede caminar sin que la viga
        pierda contacto en $A$ es
        \begin{equation*}
            x_{\max} = \frac{19}{6}\,\text{m}.
        \end{equation*}
    \end{solucion}
\end{document}